%% ============================================================
%%  6.  RESULTS
%% ============================================================
\section{Results}
\label{sec:results}

\subsection{Accuracy and Hypothesis Preservation}

Table~\ref{tab:main} summarises accuracy, Premature Elimination Rate (PER),
and Consensus Reliability across all six conditions.

\begin{table}[t]
\centering
\caption{Main results on 200 MMLU-Pro questions.
$\Delta\text{Acc}$ is absolute improvement over baseline.
PER = premature elimination rate (lower is better).
CR = consensus reliability.
Best result per column in \textbf{bold}.
Significance vs.\ baseline (McNemar): $^{*}p\!<\!0.05$, $^{**}p\!<\!0.01$.}
\label{tab:main}
\resizebox{\columnwidth}{!}{%
\begin{tabular}{lcccc}
\toprule
\textbf{Method} & \textbf{Acc} & $\Delta$\textbf{Acc} & \textbf{PER} & \textbf{CR} \\
\midrule
Baseline (Exp1)           & 54.5\% & ---        & 8.20\% & 57.1\% \\
Loser Resist.--Support (Exp2a) & 55.5\% & +1.0pp  & 9.02\% & 56.3\% \\
Loser Resist.--Metric  (Exp2b) & 58.0\% & +3.5pp$^{*}$ & 7.81\% & 58.3\% \\
Static Protection      (Exp3a) & 57.5\% & +3.0pp$^{*}$ & 6.20\% & 58.1\% \\
Diag.\ Protection      (Exp3b) & \textbf{60.0\%} & \textbf{+5.5pp$^{**}$} & \textbf{3.17\%} & \textbf{60.6\%} \\
\midrule
Learned Predictor      (Exp4)  & \PLACEHOLDER{Exp4\_ACC}\% &
                                  \PLACEHOLDER{Exp4\_DELTA\_ACC}$^{\dag}$ &
                                  \PLACEHOLDER{Exp4\_PER}\% &
                                  \PLACEHOLDER{Exp4\_CR}\% \\
\bottomrule
\end{tabular}}
\end{table}

\paragraph{Overall trend.}
Accuracy improves monotonically from Exp2b through Exp3b as triggers become
more principled.
The diagnostic trigger in Exp3b achieves $+5.5$ percentage points over
baseline (54.5\%→60.0\%, McNemar $p\!<\!0.01$), and reduces PER by
61.3\% (8.20\%→3.17\%).
Consensus reliability also rises from 57.1\% to 60.6\%, meaning that when
agents do agree under Exp3b, they are more likely to be right.

\paragraph{Exp2a hurts hypothesis preservation.}
The support-triggered loser resistance (Exp2a) is the \emph{only} condition
that worsens PER relative to baseline (+9.7\%).
Adversarial incentives without a quality filter appear to entrench agents in
their current positions indiscriminately, causing incorrect minority views to
persist alongside correct ones.

\paragraph{Diagnostic trigger outperforms categorical trigger.}
Exp3b (diagnostic) substantially outperforms Exp3a (categorical) despite
sharing the same protection mechanism, same token budget, and same prompt.
The categorical trigger activates in 39.5\% of debates including many
where no genuine imbalance exists, whereas the diagnostic trigger fires in
36.5\% of debates but only when influence asymmetry or rapid convergence is
detected.
This selectivity explains the gap: Exp3a paradoxically increases mean
influence asymmetry to 0.235 (vs.\ 0.098 at baseline), suggesting that
indiscriminate minority amplification creates its own structural imbalance.

\paragraph{Learned predictor (Exp4).}
\PLACEHOLDER{Fill in when Exp4 completes. Compare Exp4 accuracy, PER, and
CR to Exp3b. Report whether predictor is more or less selective
(intervention rate) and whether precision is higher.}

%% ----------------------------------------------------------
\subsection{Process Metrics}
%% ----------------------------------------------------------

Table~\ref{tab:process} shows the six process diagnostics.

\begin{table}[t]
\centering
\caption{Process-level diagnostic metrics (mean $\pm$ std).
Higher is better for engagement, responsiveness, influence asymmetry
\emph{inverted} ($1-\hat{I}$), balance, stability, and group welfare.
Best per column in \textbf{bold}.}
\label{tab:process}
\resizebox{\columnwidth}{!}{%
\begin{tabular}{lcccccc}
\toprule
\textbf{Method} & \textbf{Eng.} & \textbf{Resp.} & \textbf{IA-inv.} & \textbf{Bal.} & \textbf{Stab.} & \textbf{GW} \\
\midrule
Baseline (Exp1)  & .106 & .045 & .902 & .692 & \textbf{.973} & .147 \\
Loser-Support (Exp2a) & .092 & .036 & .867 & .657 & .983 & .167 \\
Loser-Metric  (Exp2b) & .105 & .046 & .880 & .696 & .978 & \textbf{.173} \\
Static Prot.  (Exp3a) & .076 & .011 & .765 & \textbf{.728} & .942 & .107 \\
Diag.\ Prot.  (Exp3b) & \textbf{.116} & \textbf{.047} & .859 & .667 & .967 & .152 \\
\midrule
Learned (Exp4)        &
  \PLACEHOLDER{E4\_eng} &
  \PLACEHOLDER{E4\_resp} &
  \PLACEHOLDER{E4\_IAinv} &
  \PLACEHOLDER{E4\_bal} &
  \PLACEHOLDER{E4\_stab} &
  \PLACEHOLDER{E4\_gw} \\
\bottomrule
\end{tabular}}
\end{table}

Exp3b achieves the highest engagement (0.116) and responsiveness (0.047),
indicating that protected agents engage more substantively rather than
passively deferring.
Exp3a shows the lowest responsiveness (0.011) and influence asymmetry
inverted (0.765), confirming that its over-broad categorical trigger
creates structural dominance rather than balanced deliberation.
Stability is highest for Exp2a (0.983) and lowest for Exp3a (0.942);
Exp3b maintains stability at 0.967 while achieving the best accuracy.

%% ----------------------------------------------------------
\subsection{Intervention Analysis}
%% ----------------------------------------------------------

\begin{table}[t]
\centering
\caption{Intervention statistics across conditions.
Turn rate = fraction of agent turns with intervention applied.
Debate rate = fraction of debates with at least one intervention.}
\label{tab:intervention}
\begin{tabular}{lcc}
\toprule
\textbf{Method} & \textbf{Turn Rate} & \textbf{Debate Rate} \\
\midrule
Loser-Support  (Exp2a) & 1.5\%  & 17.5\% \\
Loser-Metric   (Exp2b) & 7.7\%  & 30.5\% \\
Static Prot.   (Exp3a) & 6.2\%  & 39.5\% \\
Diag.\ Prot.   (Exp3b) & 7.1\%  & 36.5\% \\
Learned Pred.  (Exp4)  & \PLACEHOLDER{E4\_turnrate} & \PLACEHOLDER{E4\_debaterate} \\
\bottomrule
\end{tabular}
\end{table}

Among the diagnostic conditions, Exp3b's 7.1\% turn intervention rate
produces the best accuracy outcome, suggesting that quality matters more
than quantity of interventions.
Exp2a's very conservative 1.5\% turn rate explains why it provides the
weakest accuracy lift despite the adversarial motivation.
The learned predictor (Exp4) is expected to be more selective than Exp3b
with higher precision, as it conditions on both predicted correctness and
minimum reasoning quality.
\PLACEHOLDER{Fill in Exp4 intervention analysis once results available.}
